% Sources: [file:1][file:2]

\documentclass[12pt]{article}
\usepackage[margin=1in]{geometry}
\usepackage{amsmath,amssymb,mathtools}
\usepackage{enumitem}
\usepackage{bm}
\usepackage{hyperref}
\usepackage{fancyhdr}
\usepackage{draftwatermark}
\usepackage{xcolor}

%------------- TITLE AND METADATA -------------

\newcommand{\CourseCode}{MH1201}
\newcommand{\CourseName}{Linear Algebra II}
\newcommand{\DocType}{Tutorial 1 (Week 2) -- Problems \& Solutions}

\title{
\CourseCode\ \CourseName \\[0.3em]
\large \DocType
}
\author{
  Academic Year 2025/2026, Semester 2 \\[0.25em]
  \textit{Quantitative Research Society @NTU}
}
\date{\today}

%----------------------------------------------

%------------- HEADER & FOOTER (for page 2 onward) -------------
\fancypagestyle{main}{
  \fancyhf{}
  \fancyhead[L]{\CourseCode\ \CourseName}
  \fancyhead[R]{\href{https://qrsntu.org}{QRS@NTU}}
  \fancyfoot[C]{\thepage}
  \renewcommand{\headrulewidth}{0.4pt}
  \renewcommand{\footrulewidth}{0pt}
}

%------------- SIMPLE FOOTER (page 1 only) -------------
\fancypagestyle{plainfirst}{
  \fancyhf{}
  \renewcommand{\headrulewidth}{0pt}
  \renewcommand{\footrulewidth}{0pt}
  \fancyfoot[C]{\scriptsize\textcolor{gray}{\textit{For academic reference only. This document is provided for educational purposes and does not constitute an official question paper, answer key, or any formally authorised academic material. Its content is not guaranteed to be complete or accurate.}}}
}

%------------- WATERMARK -------------
\SetWatermarkText{Unofficial -- QRS@NTU -- qrsntu.org}
\SetWatermarkScale{1}
\SetWatermarkLightness{0.99}
%------------------------------------
% Optional: tighter list defaults (feel free to adjust)
\setlist[enumerate]{itemsep=0.3em, topsep=0.3em}
\setlist[itemize]{itemsep=0.2em, topsep=0.2em}

\setcounter{secnumdepth}{-1} % Globally removes numbers from all sections
\setcounter{tocdepth}{3}     % Ensures \subsubsection level shows in TOC



\begin{document}

%------------- COVER PAGE -------------
\maketitle
\thispagestyle{plainfirst}
\pagestyle{main}
\bigskip

%====================================================
% OVERVIEW
%====================================================
\begin{center}
  \rule{0.8\textwidth}{0.4pt}\\[4pt]
  \textbf{Overview of This Tutorial Problem Set}\\[2pt]
  \rule{0.8\textwidth}{0.4pt}
\end{center}

\noindent
This sheet focuses on core \emph{subspace} ideas and common counterexamples:
\begin{itemize}[leftmargin=*]
  \item Subspaces defined by linear constraints in \(\mathbb{R}^n\).
  \item Function spaces and subspaces defined by linear conditions (evaluation/derivative).
  \item Additive subgroups that fail scalar closure.
  \item Unions of subspaces: when they are (and are not) subspaces.
  \item Subsets of matrix spaces and quick subspace tests.
  \item Direct sums and constructing complementary subspaces in \(\mathbb{R}^n\).
\end{itemize}

\bigskip
\newpage
\tableofcontents
%====================================================
% QUESTION 1
%====================================================
\newpage
\section{Question 1 (Linear constraint in \(\mathbb{R}^3\))}
\subsection{Problem}
Determine whether
\[
\left\{(a,b,c)\in \mathbb{R}^3 \,\middle|\, a=5c\right\}
\]
is a vector subspace of \(\mathbb{R}^3\).

\subsection{Solution}

\subsubsection{Method 1: Subspace test (axioms)}
Let
\[
U \coloneqq \left\{(a,b,c)\in \mathbb{R}^3 \mid a=5c\right\}.
\]

\textbf{(i) Zero vector.} The zero vector \((0,0,0)\) satisfies \(0=5\cdot 0\), hence \((0,0,0)\in U\).

\textbf{(ii) Closure under addition.}
Take \((a_1,b_1,c_1),(a_2,b_2,c_2)\in U\). Then \(a_1=5c_1\) and \(a_2=5c_2\). Therefore
\[
a_1+a_2 = 5c_1+5c_2 = 5(c_1+c_2),
\]
so \((a_1+a_2,\; b_1+b_2,\; c_1+c_2)\in U\).

\textbf{(iii) Closure under scalar multiplication.}
Take \((a,b,c)\in U\) and \(\lambda\in \mathbb{R}\). Since \(a=5c\),
\[
\lambda a = \lambda(5c)=5(\lambda c),
\]
so \((\lambda a,\lambda b,\lambda c)\in U\).

Thus \(U\) is a subspace of \(\mathbb{R}^3\), and the final answer is
\[
\boxed{\text{\(U\) is a vector subspace of \(\mathbb{R}^3\).}}
\]

\subsubsection{Method 2: Parametrisation / spanning set}
The condition \(a=5c\) means that any vector in \(U\) has the form
\[
(a,b,c)=(5c,\; b,\; c)= b(0,1,0)+c(5,0,1).
\]
Hence
\[
U = \operatorname{span}\{(0,1,0),\ (5,0,1)\}.
\]
A span of vectors is always a subspace, so again
\[
\boxed{\text{\(U\) is a vector subspace of \(\mathbb{R}^3\).}}
\]

%====================================================
% QUESTION 2
%====================================================
\newpage
\section{Question 2 (Subspace of differentiable functions)}
\subsection{Problem}
Recall that \(\mathbb{R}^{\mathbb{R}}\) is a vector space under pointwise addition and pointwise scalar multiplication.
Show that the set of differentiable functions \(f:\mathbb{R}\to\mathbb{R}\) satisfying
\[
f'(1)=3f(2)
\]
is a vector subspace of \(\mathbb{R}^{\mathbb{R}}\).

\subsection{Solution}
Let \(\mathcal{D}\) be the set of all differentiable functions \(f:\mathbb{R}\to\mathbb{R}\), and define
\[
W \coloneqq \{ f\in \mathcal{D} \mid f'(1)=3f(2)\}.
\]
Since \(\mathcal{D}\subseteq \mathbb{R}^{\mathbb{R}}\) is itself a vector space under pointwise operations, it suffices to show \(W\) is a subspace of \(\mathcal{D}\).

\subsubsection{Method 1: Direct subspace verification}
\textbf{(i) Zero function.}
Let \(0\) denote the zero function \(0(x)=0\). Then \(0'(1)=0\) and \(3\cdot 0(2)=0\), so \(0\in W\).

\textbf{(ii) Closure under addition.}
Take \(f,g\in W\). Then \(f'(1)=3f(2)\) and \(g'(1)=3g(2)\).
Because \((f+g)'=f'+g'\), we get
\[
(f+g)'(1)=f'(1)+g'(1)=3f(2)+3g(2)=3(f(2)+g(2))=3(f+g)(2).
\]
So \(f+g\in W\).

\textbf{(iii) Closure under scalar multiplication.}
Take \(f\in W\) and \(\lambda\in\mathbb{R}\).
Because \((\lambda f)'=\lambda f'\), we get
\[
(\lambda f)'(1)=\lambda f'(1)=\lambda\cdot 3f(2)=3(\lambda f(2))=3(\lambda f)(2).
\]
So \(\lambda f\in W\).

Therefore \(W\) is a vector subspace, and the final answer is
\[
\boxed{\text{\(W\) is a vector subspace of \(\mathbb{R}^{\mathbb{R}}\).}}
\]

\subsubsection{Method 2: Kernel of a linear functional}
Define a map \(L:\mathcal{D}\to \mathbb{R}\) by
\[
L(f)\coloneqq f'(1)-3f(2).
\]
We check linearity. For \(f,g\in\mathcal{D}\) and \(\lambda\in\mathbb{R}\),
\begin{align*}
L(f+g) &= (f+g)'(1)-3(f+g)(2) \\
&= (f'(1)+g'(1)) - 3(f(2)+g(2)) \\
&= (f'(1)-3f(2)) + (g'(1)-3g(2)) \\
&= L(f)+L(g),
\end{align*}
and
\begin{align*}
L(\lambda f) &= (\lambda f)'(1)-3(\lambda f)(2) \\
&= \lambda f'(1) - 3\lambda f(2) \\
&= \lambda(f'(1)-3f(2))\\
&= \lambda L(f).
\end{align*}
Hence \(L\) is linear. Now observe:
\[
W=\{f\in\mathcal{D}\mid f'(1)=3f(2)\}=\{f\in\mathcal{D}\mid L(f)=0\}=\ker(L).
\]
The kernel of a linear map is a subspace, so
\[
\boxed{\text{\(W\) is a vector subspace of \(\mathbb{R}^{\mathbb{R}}\).}}
\]

%====================================================
% QUESTION 3
%====================================================
\newpage
\section{Question 3 (Additive closure but not scalar closure)}
\subsection{Problem}
Find a subset \(S\subseteq \mathbb{R}^2\) that satisfies the following properties:
\begin{itemize}[leftmargin=*]
  \item \(S\) is closed under vector addition and under taking additive inverses.
  \item \(S\) is not closed under scalar multiplication.
\end{itemize}

\subsection{Solution}

\subsubsection{Method 1: Integer lattice example}
Let
\[
S\coloneqq \mathbb{Z}^2=\{(m,n)\in \mathbb{R}^2\mid m,n\in\mathbb{Z}\}.
\]

\textbf{Closed under addition.} If \((m_1,n_1),(m_2,n_2)\in \mathbb{Z}^2\), then
\[
(m_1,n_1)+(m_2,n_2)=(m_1+m_2,\;n_1+n_2)\in\mathbb{Z}^2
\]
since sums of integers are integers.

\textbf{Closed under additive inverses.} If \((m,n)\in\mathbb{Z}^2\), then
\[
-(m,n)=(-m,-n)\in\mathbb{Z}^2
\]
since negatives of integers are integers.

\textbf{Not closed under scalar multiplication.} Take \((1,0)\in\mathbb{Z}^2\) and scalar \(\lambda=\tfrac12\in\mathbb{R}\). Then
\[
\tfrac12(1,0)=\left(\tfrac12,0\right)\notin \mathbb{Z}^2.
\]
Hence \(\mathbb{Z}^2\) fails scalar closure (over \(\mathbb{R}\)), so it is not a vector subspace.

Thus one valid answer is
\[
\boxed{\,S=\mathbb{Z}^2\, }.
\]

\subsubsection{Method 2: Rational lattice example}
Alternatively, let
\[
S\coloneqq \mathbb{Q}^2=\{(p,q)\in \mathbb{R}^2\mid p,q\in\mathbb{Q}\}.
\]

\textbf{Closed under addition and additive inverses.}
If \(p_1,p_2,q_1,q_2\in\mathbb{Q}\), then \(p_1+p_2\in\mathbb{Q}\), \(q_1+q_2\in\mathbb{Q}\), and \(-p_1,-q_1\in\mathbb{Q}\).
So \(\mathbb{Q}^2\) is closed under addition and taking negatives.

\textbf{Not closed under scalar multiplication by reals.}
Take \((1,0)\in\mathbb{Q}^2\) and scalar \(\lambda=\sqrt{2}\in\mathbb{R}\). Then
\[
\sqrt{2}(1,0)=(\sqrt{2},0)\notin\mathbb{Q}^2.
\]

So another valid answer is
\[
\boxed{\,S=\mathbb{Q}^2\, }.
\]

%====================================================
% QUESTION 4
%====================================================
\newpage
\section{Question 4 (Union of subspaces may fail)}
\subsection{Problem}
Find vector subspaces \(W_1\) and \(W_2\) of \(\mathbb{R}^2\) such that \(W_1\cup W_2\) is not a vector subspace of \(\mathbb{R}^2\).
\newline
\textbf{Hint:} Consider the coordinate axes of \(\mathbb{R}^2\).

\subsection{Solution}

\subsubsection{Method 1: Coordinate axes counterexample}
Let
\[
W_1=\{(x,0)\in\mathbb{R}^2\mid x\in\mathbb{R}\},\qquad
W_2=\{(0,y)\in\mathbb{R}^2\mid y\in\mathbb{R}\}.
\]
Each is a subspace (each is a line through the origin).

Now \((1,0)\in W_1\subseteq W_1\cup W_2\) and \((0,1)\in W_2\subseteq W_1\cup W_2\).
But
\[
(1,0)+(0,1)=(1,1)\notin W_1\cup W_2
\]
since \((1,1)\) lies on neither axis.
So \(W_1\cup W_2\) is not closed under addition and hence is not a subspace.

Therefore one valid choice is
\[
\boxed{\,W_1=\{(x,0)\},\ \ W_2=\{(0,y)\}\ \text{and}\ W_1\cup W_2\ \text{is not a subspace.}\,}
\]

\subsubsection{Method 2: Conceptual reason (two distinct lines)}
Let \(W_1,W_2\) be two distinct 1-dimensional subspaces (lines through the origin) in \(\mathbb{R}^2\), with \(W_1\neq W_2\).
Choose \(u\in W_1\setminus W_2\) and \(v\in W_2\setminus W_1\); such vectors exist because the subspaces are distinct.
Then \(u,v\) are linearly independent (if \(u\) were a scalar multiple of \(v\), it would lie in \(W_2\)).
Hence \(u+v\) is not in \(W_1\) (otherwise \(v=(u+v)-u\in W_1\)), and similarly \(u+v\notin W_2\).
Thus \(u,v\in W_1\cup W_2\) but \(u+v\notin W_1\cup W_2\), so the union cannot be a subspace.

Applying this to the \(x\)-axis and \(y\)-axis reproduces Method 1.

%====================================================
% QUESTION 5
%====================================================
\newpage
\section{Question 5 (When is \(W_1\cup W_2\) a subspace?)}
\subsection{Problem}
Let \(V\) be a vector space, and let \(W_1\) and \(W_2\) be vector subspaces of \(V\).
Show that \(W_1\cup W_2\) is a vector subspace of \(V\) if and only if either \(W_1\subseteq W_2\) or \(W_2\subseteq W_1\).
\newline
\textbf{Hint:} Use the Principle of Contradiction to establish the forward direction.

\subsection{Solution}
\subsubsection{Method 1: Direct contradiction (subspace-axiom argument)}
\textbf{(\(\Rightarrow\))} Assume \(W_1\cup W_2\) is a subspace of \(V\). Suppose for contradiction that neither inclusion holds.
Then there exist
\[
u\in W_1\setminus W_2,\qquad v\in W_2\setminus W_1.
\]
Since \(u,v\in W_1\cup W_2\) and \(W_1\cup W_2\) is a subspace, it is closed under addition, hence
\[
u+v\in W_1\cup W_2.
\]
If \(u+v\in W_1\), then \(v=(u+v)-u\in W_1\), contradicting \(v\notin W_1\).
If \(u+v\in W_2\), then \(u=(u+v)-v\in W_2\), contradicting \(u\notin W_2\).
Therefore the assumption was false, and thus \(W_1\subseteq W_2\) or \(W_2\subseteq W_1\).

\textbf{(\(\Leftarrow\))} If \(W_1\subseteq W_2\), then \(W_1\cup W_2=W_2\), a subspace. If \(W_2\subseteq W_1\), then
\(W_1\cup W_2=W_1\), also a subspace.

\subsubsection{Method 2: Apply the union-of-subgroups theorem}
\textbf{Theorem (Union of subgroups).} If \(H,K\) are subgroups of a group \(G\), then \(H\cup K\) is a subgroup of \(G\)
if and only if \(H\subseteq K\) or \(K\subseteq H\).

\textbf{(\(\Rightarrow\))} View \((V,+)\) as an abelian group. Since \(W_1,W_2\le (V,+)\) and \(W_1\cup W_2\) is assumed
to be a \emph{vector} subspace, it is in particular an additive subgroup of \((V,+)\). By the theorem,
\[
W_1\cup W_2 \text{ is a subgroup } \Longrightarrow \ (W_1\subseteq W_2 \ \text{or}\ W_2\subseteq W_1).
\]

\textbf{(\(\Leftarrow\))} If \(W_1\subseteq W_2\) then \(W_1\cup W_2=W_2\) is a subspace; if \(W_2\subseteq W_1\) then
\(W_1\cup W_2=W_1\) is a subspace.

\medskip
Hence,
\[
\boxed{\,W_1\cup W_2\text{ is a subspace of }V \Longleftrightarrow (W_1\subseteq W_2\ \text{or}\ W_2\subseteq W_1).\,}
\]

\newpage
\section{Question 6 (Subset of \(\mathbb{R}^{2,2}\))}
\subsection{Problem}
Determine whether
\[
\left\{
\begin{pmatrix}
a & b\\
3b-4c & a+2
\end{pmatrix}
\in \mathbb{R}^{2,2}
\ \middle|\ a,b,c\in\mathbb{R}
\right\}
\]
is a vector subspace of \(\mathbb{R}^{2,2}\).

\subsection{Solution}
Let
\[
S \coloneqq \left\{
\begin{pmatrix}
a & b\\
3b-4c & a+2
\end{pmatrix}
\ \middle|\ a,b,c\in\mathbb{R}
\right\}.
\]

\subsubsection{Method 1: Check whether the zero matrix is included}
If \(S\) were a subspace of \(\mathbb{R}^{2,2}\), it must contain the zero matrix
\[
\mathbf{0}=
\begin{pmatrix}
0&0\\0&0
\end{pmatrix}.
\]
Suppose \(\mathbf{0}\in S\). Then there exist \(a,b,c\in\mathbb{R}\) such that
\[
\begin{pmatrix}
a & b\\
3b-4c & a+2
\end{pmatrix}
=
\begin{pmatrix}
0 & 0\\
0 & 0
\end{pmatrix}.
\]
Equating entries gives
\[
a=0,\quad b=0,\quad 3b-4c=0,\quad a+2=0.
\]
But \(a=0\) and \(a+2=0\) cannot both hold, since \(a+2=0\) forces \(a=-2\).
This contradiction shows \(\mathbf{0}\notin S\).
Therefore \(S\) is not a subspace.

Hence the final answer is
\[
\boxed{\text{\(S\) is not a vector subspace of \(\mathbb{R}^{2,2}\).}}
\]



%====================================================
% QUESTION 7
%====================================================
\newpage
\section{Question 7 (Direct sum complement in \(\mathbb{R}^5\))}
\subsection{Problem}
Define a vector subspace \(V\) of \(\mathbb{R}^5\) as follows:
\[
V \coloneqq \left\{(a,b,a+b,a-b,2a)\in\mathbb{R}^5 \ \middle|\ a,b\in\mathbb{R}\right\}.
\]
Find a vector subspace \(W\) of \(\mathbb{R}^5\) such that \(\mathbb{R}^5 = V \oplus W\).

\subsection{Solution}

\subsubsection{Method 1: Basis extension via an explicit direct-sum decomposition}
Write each \(x\in V\) as
\[
(a,b,a+b,a-b,2a)=a(1,0,1,1,2)+b(0,1,1,-1,0).
\]
Hence
\[
V=\operatorname{span}\{v_1,v_2\},\qquad 
v_1=(1,0,1,1,2),\ \ v_2=(0,1,1,-1,0),
\]
so \(\dim V\le 2\), and in fact \(v_1,v_2\) are linearly independent (their first two coordinates form the standard basis of \(\mathbb{R}^2\)), hence \(\dim V=2\).

Define
\[
W \coloneqq \operatorname{span}\{e_3,e_4,e_5\}
= \{(0,0,s,t,u)\in \mathbb{R}^5 \mid s,t,u\in\mathbb{R}\}.
\]

\textbf{Claim 1: \(V\cap W=\{\mathbf{0}\}\).}
Let \(x\in V\cap W\). Since \(x\in W\), we have \(x_1=x_2=0\). Since \(x\in V\), there exist \(a,b\in\mathbb{R}\) such that
\[
x=(a,b,a+b,a-b,2a).
\]
Comparing coordinates gives \(a=x_1=0\) and \(b=x_2=0\), hence \(x=\mathbf{0}\). Therefore \(V\cap W=\{\mathbf{0}\}\).

\textbf{Claim 2: \(V+W=\mathbb{R}^5\).}
Let \(x=(x_1,x_2,x_3,x_4,x_5)\in\mathbb{R}^5\). Set
\[
v \coloneqq (x_1,x_2,x_1+x_2,x_1-x_2,2x_1)\in V,
\]
and define \(w\coloneqq x-v\). Then
\[
w=(0,0,\ x_3-(x_1+x_2),\ x_4-(x_1-x_2),\ x_5-2x_1)\in W.
\]
Thus \(x=v+w\) with \(v\in V\) and \(w\in W\), proving \(V+W=\mathbb{R}^5\).

Since \(V+W=\mathbb{R}^5\) and \(V\cap W=\{\mathbf{0}\}\), it follows that
\[
\mathbb{R}^5 = V \oplus W.
\]
Hence a valid choice is
\[
\boxed{\,W=\{(0,0,s,t,u)\in\mathbb{R}^5:\ s,t,u\in\mathbb{R}\}=\operatorname{span}\{e_3,e_4,e_5\}.\,}
\]

\subsubsection{Method 2: Kernel--image decomposition using a linear map (rank--nullity)}
Define a linear map \(T:\mathbb{R}^5\to\mathbb{R}^2\) by
\[
T(x_1,x_2,x_3,x_4,x_5)\coloneqq (x_1,x_2).
\]
Then \(T\) is surjective, hence \(\operatorname{rank}(T)=2\). Moreover,
\[
\ker(T)=\{(0,0,s,t,u)\in\mathbb{R}^5:\ s,t,u\in\mathbb{R}\}.
\]
We claim that \(T|_V:V\to\mathbb{R}^2\) is an isomorphism. Indeed, for any \((a,b)\in\mathbb{R}^2\),
\[
T(a,b,a+b,a-b,2a)=(a,b),
\]
so \(T(V)=\mathbb{R}^2\) (surjectivity). Also, if \(x\in V\) and \(T(x)=(0,0)\), then \(x\) has the form
\[
x=(a,b,a+b,a-b,2a)\in V\quad\text{with }(a,b)=(0,0),
\]
hence \(x=\mathbf{0}\), so \(T|_V\) is injective.

Now let \(W\coloneqq \ker(T)\). For any \(x\in\mathbb{R}^5\), define \((a,b)\coloneqq T(x)=(x_1,x_2)\), and let
\[
v\coloneqq (a,b,a+b,a-b,2a)\in V.
\]
Then \(T(x-v)=T(x)-T(v)=(a,b)-(a,b)=(0,0)\), hence \(x-v\in \ker(T)=W\). Therefore \(x=v+w\) for some \(v\in V\), \(w\in W\), so \(\mathbb{R}^5=V+W\).

Finally, if \(x\in V\cap W\), then \(x\in V\) and \(x\in \ker(T)\) implies \(T(x)=(0,0)\). Since \(T|_V\) is injective, this forces \(x=\mathbf{0}\). Thus \(V\cap W=\{\mathbf{0}\}\), and hence \(\mathbb{R}^5=V\oplus W\).

Consequently,
\[
\boxed{\,W=\ker(T)=\operatorname{span}\{e_3,e_4,e_5\}.\,}
\]
\end{document}
